\chapter{Software architecture}\label{ch:arch}

This chapter describes how the code is organised, so that the algorithms of the preceding chapters can be found, tested
and changed.

\section{Two layers}
The package \code{rs3d} has two layers with a strict dependency direction.

\begin{figure}[tbp]
\centering
\begin{tikzpicture}[node distance=5mm]
  \node[boxblue,minimum width=0.9\linewidth] (app) {\textbf{app/} --- NiceGUI pages, ECharts option builders, strings, state\\
    {\small status, live, helicorder, data, settings, guide \textbullet\ charts \textbullet\ detail panel}};
  \node[box,minimum width=0.9\linewidth,below=of app] (core) {\textbf{core/} --- signal processing and data access; no user-interface imports\\
    {\small spectro, helicorder, health, coverage, export, live, store, archive, seedlink, response, status}};
  \node[boxdark,minimum width=0.9\linewidth,below=of core] (lib) {ObsPy \textbullet\ SciPy \textbullet\ NumPy \textbullet\ matplotlib \textbullet\ paramiko \textbullet\ keyring \textbullet\ psutil};
  \draw[flow] (app) -- (core);
  \draw[flow] (core) -- (lib);
\end{tikzpicture}
\caption{Layers of the application. Arrows point from a layer to the layer it uses.}\label{fig:layers}
\end{figure}

\begin{description}[leftmargin=2.2cm,style=nextline]
\item[\code{core/}] every algorithm in this manual. These modules know nothing of the browser; they take and return ObsPy
  streams, NumPy arrays and plain data classes, which makes them testable in milliseconds with synthetic signals.
\item[\code{app/}] presentation. The chart builders in \file{app/charts.py} are \emph{pure functions} that return the
  option dictionaries of the charting library, so even the shape of a plot is unit-tested without a browser.
\end{description}

\section{Where each algorithm lives}
\Cref{tab:modules} maps every topic of this manual to the module and the main functions that implement it.
\begin{table}[tbp]
\centering\footnotesize
\caption{Algorithm $\to$ module.}\label{tab:modules}
\begin{tabular}{@{}L{0.20\linewidth}L{0.25\linewidth}L{0.40\linewidth}@{}}
\toprule
Topic & Module & Main functions\\ \midrule
Day-file layout, import & \file{core/store.py} & \code{day\_path}, \code{load}, \code{import\_file}\\
Response, units, filtering & \file{core/response.py}, \file{export.py} & \code{get\_inventory}, \code{process}\\
Spectra, spectrogram, noise & \file{core/spectro.py} & \code{amplitude\_spectrum}, \code{spectrogram}, \code{window\_for}, \code{psd\_summary}\\
Completeness, coverage & \file{core/health.py}, \file{coverage.py} & \code{health}, \code{coverage}\\
Helicorder rows and windows & \file{core/helicorder.py} & \code{build\_rows}, \code{scale\_for}, \code{click\_to\_time}, \code{zoom\_window}, \code{box\_window}\\
Decimation, live buffer & \file{core/live.py} & \code{minmax\_decimate}, \code{RingBuffer}, \code{LiveStream}\\
Shake access & \file{core/seedlink.py}, \file{shell.py}, \file{archive.py} & \code{probe}, \code{SeedLinkSource}, \code{sync}\\
Plot shapes & \file{app/charts.py} & \code{helicorder\_options}, \code{spectrogram\_options}, \code{psd\_options}\\
Detail view, clipboard & \file{app/detail.py}, \file{core/export.py} & \code{DetailPanel}, \code{write\_clipboard\_png}\\
\bottomrule
\end{tabular}
\end{table}

\section{Reproducibility by construction}
The Data page can show, and save, a short Python program that reproduces exactly the processing behind a preview: it
reads the same day files, applies the same response removal and the same band-pass, and prints the stream. The program is
generated by \file{core/codegen.py} from the same \code{Selection} object that the application itself uses, so the code
and the screen cannot drift apart. A researcher can therefore take any number off the screen and recompute it outside the
application.

\section{Testing}
Three levels of test protect the algorithms.
\begin{enumerate}
\item \textbf{Unit tests} of \file{core/} with synthetic data: a sine must peak in its own bin and read its own amplitude, a
  spike must survive decimation, a gap must stay a gap, a click must map to its row and time.
\item \textbf{Page tests} with NiceGUI's user simulation: each page renders with a fake device; controls change state as
  expected.
\item \textbf{Real-browser tests} (Playwright with Google Chrome, marker \code{e2e}): import a file, open the helicorder, click
  a row, zoom with a dragged rectangle, open the clipboard and download the picture. Each such test has been shown to fail when
  the behaviour it protects is deliberately broken.
\end{enumerate}
A fourth marker, \code{device}, selects tests that need a real Shake and is skipped by default.

\section{Configuration and storage}
All files live in one folder (on Linux \file{\textasciitilde/.local/share/rs3d}, overridable with the
\code{RS3D\_DATA\_DIR} environment variable): \file{config.toml} (devices), \file{settings.toml} (time zone, last station,
units and band-pass), \file{known\_hosts}, \file{sds/} (data), \file{response/} (StationXML cache), \file{logs/}. Passwords
are kept in the operating system's key store through \code{keyring}, never in these files.
